\documentclass[12pt]{article}
\usepackage[utf8]{inputenc}
\usepackage[spanish]{babel}
\usepackage[a4paper, margin=2cm]{geometry}
\usepackage{tikz}
\usetikzlibrary{babel}
\usepackage[most]{tcolorbox}
\usepackage{amsmath}

\renewcommand{\familydefault}{\sfdefault}

% Definición de colores
\definecolor{medblue}{RGB}{43, 108, 176}
\definecolor{medteal}{RGB}{49, 151, 149}
\definecolor{medlight}{RGB}{235, 248, 255}
\definecolor{meddark}{RGB}{26, 32, 44}
\definecolor{iron}{RGB}{107, 114, 128}      % Gris para la pesa
\definecolor{rope}{RGB}{217, 119, 6}        % Color cuerda
\definecolor{bedblue}{RGB}{96, 165, 250}    % Azul para la camilla

\begin{document}
\pagestyle{empty}

\begin{tcolorbox}[
    enhanced, 
    colback=medlight!30, 
    colframe=medblue, 
    title={\Large \textbf{Trabajo y Energía: Rehabilitación Física}}, 
    halign title=center, 
    fonttitle=\bfseries, 
    arc=4mm, 
    boxrule=1.5pt,
    shadow={2mm}{-2mm}{0mm}{black!30!white}
]

    \vspace{0.2cm}
    \begin{tcolorbox}[colback=white, colframe=medteal, arc=2mm, boxrule=1.2pt, title=\textbf{Caso Clínico / Problema}]
        Un paciente en una camilla de tracción levanta repetidamente una pesa de \textbf{5 kg} a una altura de \textbf{0.4 m} a velocidad constante. Calcule el trabajo mecánico realizado por el paciente en un solo levantamiento.
    \end{tcolorbox}

    \vspace{0.4cm}
    
    \begin{center}
    \begin{tikzpicture}[scale=1.1]
        
        % Soporte y Polea
        \filldraw[meddark!20, draw=meddark, thick, rounded corners=2pt] (1.5, 3.5) rectangle (2.5, 3.8);
        \draw[thick, meddark, line width=2pt] (2, 3.5) -- (2, 3);
        \filldraw[gray!30, draw=meddark, thick] (2, 2.7) circle (0.3cm);
        \filldraw[meddark] (2, 2.7) circle (0.05cm); % Eje de la polea
        
        % Cuerda
        \draw[ultra thick, rope] (-0.5, 1) -- (1.7, 2.7); % Hacia el paciente
        \draw[ultra thick, rope] (2.3, 2.7) -- (2.3, 0.5); % Hacia la pesa
        
        % Pierna del paciente (simplificada)
        \filldraw[bedblue!50, draw=meddark, thick, rounded corners=5pt] (-2.5, 0.5) rectangle (-0.5, 1.5);
        \node[font=\bfseries, text=meddark] at (-1.5, 1) {Paciente};
        
        % Pesa (Masa)
        \filldraw[iron!80, draw=meddark, thick, rounded corners=3pt] (1.8, -0.5) rectangle (2.8, 0.5);
        \node[font=\large\bfseries, text=white] at (2.3, 0) {5 kg};
        
        % Vectores de Fuerza y Distancia
        \draw[->, ultra thick, red!80!black] (2.3, 0.5) -- (2.3, 1.8) node[midway, left, font=\bfseries] {$F$};
        \draw[->, ultra thick, medteal] (3.3, -0.5) -- (3.3, 1.2) node[midway, right, font=\bfseries] {$h = 0.4 \text{ m}$};
        
        % Guías punteadas para la altura
        \draw[dashed, meddark!60, thick] (2.8, -0.5) -- (3.5, -0.5);
        \draw[dashed, meddark!60, thick] (2.8, 1.2) -- (3.5, 1.2);

    \end{tikzpicture}
    \end{center}

    \vspace{0.4cm}
    \begin{tcolorbox}[colback=medteal!10, colframe=medteal, arc=2mm, boxrule=1.2pt, title=\textbf{Desarrollo Matemático y Solución}]
        El trabajo mecánico ($W$) realizado para elevar un objeto a velocidad constante contra la gravedad equivale al cambio en su energía potencial gravitatoria ($W = m \cdot g \cdot h$).
        \begin{align*}
            W &= (5 \text{ kg}) \times (9.8 \text{ m/s}^2) \times (0.4 \text{ m}) \\[1ex]
            W &= 49 \text{ N} \times 0.4 \text{ m} = \mathbf{19.6 \text{ Joules}}
        \end{align*}
        \textbf{Resultado:} El paciente realiza un trabajo mecánico de \textbf{19.6 Joules (J)} en cada repetición del ejercicio.
    \end{tcolorbox}
    
\end{tcolorbox}

\end{document}